\section{Introduction}

A simple closed geodesic on the modular punctured torus carries an
integer: its hyperbolic trace divided by three. The integers obtained
in this way are the Markov numbers. Their construction remembers a
curve, a rational slope, or a position in a tree of Markov triples.
The number itself forgets this information. The Markov--Frobenius
uniqueness conjecture~\cite{Frobenius1913} asks whether that loss of information creates
any ambiguity beyond the symmetries already present in the construction.
In its geometric form, it asks whether two simple closed geodesics of
equal length on the modular torus must be related by an isometry.

This paper studies the question with a parameter attached to the
\emph{pair} of curves. Choose one curve as the first element $A$ of a
simple basis $(A,B)$. The second curve then has primitive homology
$(h,N)$. For a normalized representative $2\leqslant h<N$, write $N=hk+t$ with
$k\geqslant1$ and $0<t<h$, and a Christoffel word represents it as
$h$ occurrences of $A$ separated by gaps $B^k$ and $B^{k+1}$.
We fix the occurrence count $h$ and allow both the gap quotient $k$
and the integral character of the basis to vary. The question is whether
the trace of this word can equal the trace of $A$. The endpoint cases
$N=1$ and $h=1$ are covered separately by a reflection argument.

There are two obstacles to making this a finite problem. First, a
basis with very large traces can still represent a curve of small trace
when it is far from the minimum of its Markov ray. Second, equality
of traces must be separated from equality explained by an isometry of
the surface. An arbitrary change of homology basis answers neither
question. Our argument centers each ray at its actual integral minimum,
bounds all remaining indices and parameters, and retains the marking
throughout the resulting finite calculation.

\subsection{Main conclusions}

Our main quantitative result is Theorem~\ref{thm:finite-decision}.
For each $h\geqslant2$, every primitive mechanical trace equality in
the preceding family lies in a fully explicit domain with $O(h^5)$
exact comparisons on $O(h^3)$-bit integers. The isometry status of
every equality is decided by twelve fixed integral actions. This is a
polynomial bound in the numerical value of $h$, rather than in the
length of its binary encoding.

The two parameter estimates behind this result are particularly simple.
Write $p=\tr B/3$ after centering the ray. If the base index $l$ of
the product is positive, then
$$
 p<4^{h+1}.
$$
If $l=0$, mechanical balance gives
$$
 p<(100h^2)^{h+1}.
$$
These are analytic bounds, valid over all positive integral Markov
characters. Although their heights
are large, the number of Markov occurrences below them is small enough
for a polynomial decision bound. Counting occurrences instead of scanning
integer labels is an essential part of the argument.

The finite applications give two complete rigidity statements:
all primitive occurrence families $2\leqslant h\leqslant6$ contain
only isometric equalities, and every equal-length simple pair with
intersection number at most $24$ is isometric. No upper bound on the common trace is assumed. Both proofs combine the analytic
reductions with an exact finite computation, together with separate
arguments for the remaining intersection-$19$ and intersection-$23$
types; the computation is
described where it is used, and the data it produces accompany the
article as supplementary material. Section~\ref{sec:rigidity} specifies the geometric normalization
and the precise scope of these conclusions.

\subsection{Why a comparison with commuting traces helps}

The mechanism is an integrality contradiction. If $s>1$, the commuting
baselines
$$
 B_i=\frac{s^i+s^{-i}}3
$$
satisfy $3B_aB_b-B_{a+b}=B_{|a-b|}$. Their indices therefore follow
the Euclidean algorithm under Vieta mutation. When the indices reach
zero, the corresponding baseline is $B_0=2/3$.
Actual positive integral Markov triples obey the same polynomial
mutation, but lie on a different level of the trace cubic. If their
errors relative to the baselines remain sufficiently small, the final
actual entry would be a positive integer below one.

The comparison has a nonzero forcing term, so small initial errors
alone do not suffice. At positive product index, the quotient form of
Vieta mutation keeps a weighted error bounded by the same constant at
every step. At minimum index, cyclic balance first bounds the mixed
matrix paths. The equality then improves the relevant block exponents
to positive multiples of a greatest common divisor. Passing to spectral
coordinates turns a mutation into subtraction with a small additive
error, whose propagation is controlled by the integer rows of the
Euclidean algorithm. This is the reason for using two complementary
error estimates.

\subsection{Relation to earlier work}
\label{sec:literature}

The correspondence between Markov arithmetic and the geometry of the
punctured torus is classical. Cohn related Markov forms to simple
geodesics~\cite{Cohn1971}, and Reutenauer gave a constructive bijection
between proper Christoffel words and proper Markov
triples~\cite[Theorem~1]{Reutenauer2009}. These correspondences retain a
whole triple or a marked curve. They do not assert that its largest
coordinate, or the trace of the curve alone, determines it. We use the
classical dictionaries to enumerate marked occurrences and retain their
homology throughout the calculation.

The parameter in this paper is attached to a pair of curves. After choosing
one member as a basis element, the other has primitive homology
$(h,hk+t)$ and a prescribed Christoffel gap pattern. This differs from
fixing a numerator in one global Farey parametrization. The ordering
theorems of Lagisquet, Pelantov{\'a}, Tavenas and Vuillon~\cite{LagisquetEtAl2021}
and of Lee, Li, Rabideau and Schiffler~\cite{LeeEtAl2023} compare Markov
values along specified directions in that global parametrization.
Our bounds allow the integral character of the chosen basis to vary.
Changing the basis transforms both curves and their character; it cannot
be used to reset the first curve while leaving its trace equation unchanged.

Effective finiteness on character varieties is also established. Whang
proves effective bounds for integral points on geometrically irreducible
nonintegrable curves in relative character varieties, and describes the
multitwist exceptions in the more general nondegenerate
case~\cite[Theorems~1.2--1.3]{Whang2020}. Applied to a fixed trace-equality
curve, this method requires verification of the exceptional hypotheses
component by component. The present argument gives explicit bounds in
the pair parameter $h$ and a direct finite comparison domain, without
computing those algebraic curves. Its quantitative conclusion is
$O(h^5)$ trace comparisons on $O(h^3)$-bit integers. This is a bound in
the numerical value of $h$, and does not bound $h$ itself.

The geometric part uses equally classical structure. The length of a
primitive simple class on a punctured torus is determined by its primitive
homology vector, and these lengths extend to a norm~\cite{McShaneRivin1995}.
McShane and Parlier prove that, with the boundary length fixed, the locus
where two distinct simple curves have equal length is an embedded
path~\cite[Theorem~1.3]{McShaneParlier2008}. We use this theorem
for the reflection cases. The remaining arithmetic classifications retain
the full marking and test the actual finite isometry group of the modular
torus. Homological equivalence under an arbitrary change of basis is not
an isometry test at a fixed metric.

The comparison model for the analytic estimates has a useful interpretation
on the Cayley cubic. For full traces, put
$$
F(X,Y,Z)=X^2+Y^2+Z^2-XYZ.
$$
The Fricke identity gives $F=\operatorname{tr}[A,B]+2$, so the cusped
Markov character lies on $F=0$~\cite{Goldman2003}. Commuting diagonal
matrices instead give
$$
(X,Y,Z)=(u+u^{-1},v+v^{-1},uv+(uv)^{-1}),\qquad F=4.
$$
The torus cover and monomial action on this second level are
classical~\cite[Sections~1.5 and~2.1]{CantatLoray2009}. Taking
$(u,v)=(s^a,s^b)$ turns a Vieta move into subtraction of the exponents.
The same polynomial move preserves the actual Markov level separately.
Our estimates control the discrepancy along these two evolutions;
after division of the traces by three, their level difference is $4/9$.
There is no identification of the two character surfaces in this argument.

Positive trace expansions and quadratic norm gaps have substantial prior
uses. Fisac's necklace formula already records the dependence of a trace
on cyclic runs~\cite[Theorem~1.1 and Section~4]{Fisac2025}.
Altassan and Luca use Binet expansions and quadratic norms to bound
solutions of Markov-type equations in Lucas
sequences~\cite[Theorem~1.1 and Lemmas~2.1--2.2]{AltassanLuca2021}.
Their sequences have specified initial values. Here a centered Markov ray
has coefficients satisfying $cd=p^2/(9p^2-4)$, and these coefficients
vary with the integral character. The needed estimates therefore retain
both coefficients, all mixed matrix paths and the actual Vieta mutations.
At minimum index, mechanical balance controls the mixed paths, while
spectral-angle errors propagate through the Euclidean algorithm with
polynomial amplification. The contradiction is an actual positive
integer forced below one.

Finally, arithmetic restrictions on the largest Markov entry, such as
Button's ideal-theoretic criteria~\cite{Button2001}, organize partial
uniqueness by a different parameter. The applications here concern complete
families of relative curve positions: occurrence counts two through six,
and all intersections at most twenty-four. No upper bound on the common
Markov value is assumed in these classifications. These classifications and
the effective decision procedure leave open whether a nonisometric equality
occurs at some larger occurrence count.


\subsection{Organization}

Section~\ref{sec:rays} fixes the curve, word and ray conventions.
Section~\ref{sec:gap} bounds the product index and the original gap
quotient using positive paths and a quadratic norm.
Sections~\ref{sec:positive} and~\ref{sec:spectral} prove the two
parameter bounds. Section~\ref{sec:decision} gives the complete finite
domain and its isometry test. Section~\ref{sec:rigidity} applies the
procedure to the stated families; the exceptional intersection argument
and the description of the finite domains are in the appendix.
The final section explains which part of the uniqueness problem remains.
